\chapter{Building a specification}\label{ch:spec}
\chapterintro{The specification is the one document that says what problem is solved. This chapter builds one by hand
in the workspace, a small steel bracket with a solid pad and a bolt hole, and then explains each part of the document:
the coordinate frame, units, material, supports, loads, keep-in and keep-out regions, the volume fraction and the
mesh.}

\section{Tutorial: a cantilever bracket}\label{sec:tutorial}
The bracket is a 100 x 50 x 50~mm steel block, fixed at one end and loaded at the other, with a mounting pad that must
stay solid and a hole that must stay empty. The screenshots were recorded with the ladder capped at
$1.2\times10^{7}$ elements, so that the release solves in about a minute.

\subsection{Start a new design}
Choose \gui{New design\ldots} in the problem picker. You get an empty 100 x 50 x 50~mm box in steel ($E = 210000$,
$\nu = 0.3$, in N and mm) with a volume fraction of 0.3 and 20 elements along the shortest edge. The \gui{Checks} panel
says what is missing, and the ladder stays hidden until the specification is valid (Figure~\ref{fig:newdesign}).

\begin{figure}[t]
\wideshot{02_new_design.png}
\caption{A new design. The checks say what is missing and how to add it.}
\label{fig:newdesign}
\end{figure}

Select \gui{Domain} in the tree to set the size, the units, the volume fraction (\gui{volume \%}), the mesh (by element
size $h$ or by cells along the shortest edge) and an optional minimum feature size. Select \gui{Material} to set its
name, Young's modulus $E$ and Poisson's ratio $\nu$.

\subsection{Add a support}
Press \keycap{S} (or click \gui{Support}) to arm the support tool. The toolbar asks you to click a face of the domain
(Figure~\ref{fig:supporttool}). Click the face at $x = 0$. It becomes a support fixed in all three directions, the tool
returns to \gui{Select}, and the new support is selected in the inspector. There you can change the region kind (face,
patch, edge, box or cylinder) and untick fixed directions to make a roller.

\begin{figure}[t]
\wideshot{03_support_tool.png}
\caption{The support tool armed: click the face to fix.}
\label{fig:supporttool}
\end{figure}

\subsection{Add a load}
Press \keycap{L} and click the face at $x = 100$. A 1~kN load appears on a small patch centred where you clicked,
pointing into the face. Edit it in the inspector: here the force becomes 2000~N downward (force $y = -2000$) and the
patch 10 x 20~mm (Figure~\ref{fig:loadinspector}). Each edit takes effect as you type and can be undone.

\begin{figure}[t]
\wideshot{04_load_inspector.png}
\caption{A load placed by clicking and then edited in the inspector.}
\label{fig:loadinspector}
\end{figure}

\subsection{When a check fails}
The rules run \CheckDelay~ms after every edit. Untick $y$ and $z$ on the support, so that it fixes only $x$: rule
\rref{R003} reports the rigid-body motions left free (translation in $y$ and $z$, rotation about $x$), the support gets
an \gui{error} badge in the tree, and the ladder locks (Figure~\ref{fig:ruleerror}). Tick $y$ and $z$ again and the error
clears. Errors block solving; warnings and notes do not.

\begin{figure}[t]
\centering
\shot[0.30\linewidth]{06_rule_error.png}
\caption{Rule R003: the support leaves rigid-body motion free. The ladder stays locked until the support is fixed.}
\label{fig:ruleerror}
\end{figure}

\subsection{Keep-in and keep-out regions}
Press \keycap{K} and click a face to add a keep-in region (always solid), or \keycap{O} for a keep-out region (always
empty). Each starts as a slab 10\% thick against that face. In the inspector a keep-in or keep-out region can be a box
or a cylinder, and you can switch it between keep-in and keep-out.

In Figure~\ref{fig:keep} the keep-in region was resized to an 8 x 20 x 30~mm pad under the load, using its min/max
fields. The keep-out region was converted to a cylinder of radius 8~mm along $z$, centred at $x = 40$, which leaves a
through-hole for a bolt or a cable.

\begin{figure}[t]
\wideshot{05_keepin_keepout.png}
\caption{The finished setup: support (blue), load on a solid pad (green) and a keep-out cylinder (pink). All rules
pass.}
\label{fig:keep}
\end{figure}

\subsection{Undo and redo}
\keycap{Ctrl}+\keycap{Z} undoes and \keycap{Ctrl}+\keycap{Shift}+\keycap{Z} or \keycap{Ctrl}+\keycap{Y} redoes, across
your edits and agent edits alike. Every step is a record in the timeline. Chapter~\ref{ch:ladder} continues with this
bracket and solves it.

\section{The document}\label{sec:document}
\index{specification!document}
Whatever you do in the workspace ends up as one JSON document. This is the reference specification of suite task b01:

\begin{jsoncode}
{
  "schema_version": "agentto.spec/v0",
  "name": "b01 cantilever block",
  "units": {"length": "mm", "force": "N"},
  "domain": {"size": [200, 100, 100]},
  "material": {"name": "steel", "E": 210000.0, "nu": 0.3},
  "volume_fraction": 0.3,
  "supports": [{"id": "wall", "region": {"kind": "face", "face": "x-"}}],
  "loads": [{"id": "pad", "force": [0, -1000, 0],
             "region": {"kind": "patch", "face": "x+", "center": [50, 50],
                        "size": [10, 10], "shape": "rect"}}],
  "passive": [],
  "resolution": {"h": 1.0}
}
\end{jsoncode}

The schema is strict: unknown fields are rejected, every \code{id} must be unique across supports, loads and passive
regions, and a document needs at least one support, one load, a volume fraction and a resolution. A draft that is not
complete yet (no load, say) is kept as a draft; the tools report what is missing and the rules do not run until it is
complete. Appendix~\ref{app:spec} lists every field.

\subsection{Coordinate frame}
\index{coordinate frame}
The domain is an axis-aligned box from $(0,0,0)$ to the size you give. All coordinates are absolute, in that frame. The
faces are named \code{x-}, \code{x+}, \code{y-}, \code{y+}, \code{z-} and \code{z+}: \code{x-} is the face at $x = 0$,
\code{x+} the face at $x = L_x$. Which axis is ``up'' is yours to choose; the brief b06, for example, says ``y up''.

A patch on a face gives its centre and full size in the face's two in-plane axes, always in increasing axis order:
$(y,z)$ on $x$ faces, $(x,z)$ on $y$ faces, $(x,y)$ on $z$ faces. This is the most common source of misplaced regions.
\tool{list_regions} prints the in-plane axes of every face, and the viewport shows the result.

\subsection{Units}
\index{units}
Lengths are in mm, m or in; forces in N, kN or lbf. The units label the numbers; they do not convert them.

\begin{warnbox}[Units relabel; they do not convert]
Switching from mm to m changes the label on the numbers, not the numbers. Enter every value in the units shown, and
keep $E$ consistent with them: in N and mm, $E$ is in N/mm$^2$ (MPa), so steel is 210000. In kN and mm, steel is 210. In
the paper's second-provider benchmark, two of the four remaining wrong answers were exactly this slip: the agent worked
in kN and mm but left $E$ in MPa, a material 1000 times too stiff.
\end{warnbox}

\subsection{Material}
An isotropic linear-elastic material: a name, $E$ (positive) and $\nu$ (between $-1$ and $0.5$). The name is a label;
only $E$ and $\nu$ enter the computation. Because the objective is compliance under one load case, the optimal layout
does not depend on $E$, but the reported compliance does, and the specification must state the problem the engineer
meant.

\subsection{Regions}\label{sec:regions}
\index{regions}
Supports and loads select grid \emph{nodes}; keep-in and keep-out regions select \emph{elements} by their centres.
Table~\ref{tab:regions} lists the five region kinds.

\begin{table}[h]
\caption{Region kinds. All lengths are in the spec's length unit; ``in-plane order'' is the increasing axis order of
the face's two in-plane axes.}\label{tab:regions}
\begin{fixtable}{@{}l>{\raggedright\arraybackslash}X>{\raggedright\arraybackslash}p{0.26\linewidth}@{}}
Kind & Fields & Valid for \\ \midrule
\code{face} & \code{face}: one of \code{x-}, \code{x+}, \code{y-}, \code{y+}, \code{z-}, \code{z+} & supports, loads \\
\code{patch} & \code{face}; \code{center} and \code{size} in in-plane order; \code{shape} \code{rect} (default) or
\code{ellipse} (the inscribed ellipse, e.g. a washer) & supports, loads \\
\code{edge} & \code{axis}; \code{at}: the other two coordinates in increasing order; optional \code{span} [start, end]
along the axis & supports, loads \\
\code{box} & \code{min}, \code{max}: opposite corners, \code{min} $<$ \code{max} on every axis & supports, loads,
keep-in/out \\
\code{cylinder} & \code{axis}; \code{center}: the axis position in the other two coordinates; \code{radius};
\code{span} [start, end] along the axis & supports, loads, keep-in/out \\
\end{fixtable}
\end{table}

Selection uses a rounding tolerance of $10^{-9}$ element edges. A patch with no node within half its size of its
centre, or an edge between grid lines, is snapped to the nearest grid line and reported by rule \rref{R010}; nothing is
snapped silently. A region that selects nothing is an error (\rref{R002} for a support, \rref{R004} for a load) or a
warning (\rref{R011} for a passive region).

\subsection{Supports}
A support fixes the nodes of its region in the directions listed in \code{fix}; the default is all three. A roller
fixes fewer: \code{"fix": ["y"]} on a bottom edge leaves the span and width directions free. Supports only need to
remove the six rigid-body modes together; rule \rref{R003} checks exactly that.

\subsection{Loads}
A load is a total force vector, in the spec's force unit, shared equally by the nodes of its region. This keeps the
total force the same on every mesh, but it is not a pressure integration: on a curved or partial region the nodal
shares are equal, not area-weighted. All loads act together as one load case. A load component that falls on a fixed
degree of freedom is dropped, because it would go straight into the support, and rule \rref{R005} reports the share.

\subsection{Keep-in and keep-out regions}
\index{keep-in region}\index{keep-out region}\index{passive region|see{keep-in region, keep-out region}}
In the document these are \code{passive} items with \code{kind} \code{solid} (keep-in) or \code{void} (keep-out), each
a box or a cylinder. An element whose centre falls inside is held at density 1 or 0 at every iteration and is not a
design variable. Keep-out elements keep a small stiffness rather than zero, so their nodes stay well defined. The
certificate checks that every keep-in element ends at exactly 1 and every keep-out element at exactly 0.

\subsection{Volume fraction}\label{sec:vf}
\index{volume fraction}
The volume fraction $v^*$ is the share of the \emph{design} volume, which is the domain minus every keep-in and keep-out
element:
\begin{equation}
  \frac{1}{|\mathcal D|}\sum_{e\in\mathcal D}\bar\rho_e \le v^*,\qquad \bar\rho_e = 1\ (e\in\mathcal S),\quad
  \bar\rho_e = 0\ (e\in\mathcal V),
\end{equation}
where $\mathcal D$ are the design elements, $\mathcal S$ the keep-in and $\mathcal V$ the keep-out elements, and
$\bar\rho$ the filtered and projected density. Adding a keep-in pad therefore adds material on top of the budget; it
changes both the feasible set and the total material. That is why an agent that invents a pad the brief did not ask for
has changed the problem, even if the result looks stiffer.

\subsection{Resolution and the minimum feature}
\index{resolution}
Give the mesh as an element edge \code{h} or as \code{cells\_short\_edge}, the number of elements along the shortest
side; exactly one of the two. The edge must divide all three sides into whole numbers of elements, or rule \rref{R001}
fires. This resolution is the \emph{base} mesh from which the ladder is built. The release rung is the finest rung that
fits on the GPU and is often finer: for b06 the spec says 1~mm and the release rung is 0.25~mm
(Section~\ref{sec:rungs}).

The optional \code{min\_feature} is the smallest member you need. It does not constrain the optimizer; it only feeds
rule \rref{R009}, which warns when that size spans fewer than about 2.4 elements and would be filtered out.

\section{Editing through the tools}\label{sec:editing}
Every edit in the workspace is a call to one of the editing tools, logged like an agent's call. The same calls work from
Python and over MCP:

\begin{itemize}
  \item \tool{new_spec} starts a box domain; \tool{set_spec} replaces the whole document.
  \item \tool{set_params} sets scalar fields (name, size, units, $E$, $\nu$, material name, volume fraction, $h$ or
  \code{cells\_short\_edge}, \code{min\_feature}); only the arguments given change.
  \item \tool{add_support}, \tool{add_load} and \tool{add_passive} add items; \tool{update_item} changes fields of one;
  \tool{remove_item} deletes one.
  \item \tool{undo} and \tool{redo} step through the history.
\end{itemize}

A failed edit leaves the specification unchanged. Every successful edit is undoable and records the hash of the
resulting specification in the log, so any later result can be tied to the exact document it was solved from.

\begin{tipbox}[Ground before you solve]
\tool{list_regions} reports, for every item, its physical bounding box and how many nodes or elements it selects at a
given mesh, with any snapping. \tool{explain_spec} prints the whole specification in plain sentences, and
\tool{render} draws the supports, loads and passive regions as a picture. All three are cheap; use them before the
first solve.
\end{tipbox}
